\begin{frame}[t]{Explain the full workflow in sixty seconds}

\textbf{Awareness:} establish the current case and what is known.\par\medskip
\textbf{Planning:} choose an appropriate numerical method and form candidates.\par\medskip
\textbf{Simulation:} calculate each candidate's consequences in a declared model.\par\medskip
\textbf{Verification:} check physical limits, identity, evidence and permissions.\par\medskip
\textbf{LLM optimizer:} propose a bounded next trial, then return through simulation and verification.\par\medskip
\textbf{Operator review:} inspect the recommendation and unresolved assumptions before any authorized action.
\par\medskip
\textcolor{forest}{\textbf{Partner exercise:}} use one of today's cases and name the actual method and output at every stage.

\slidetransition{Which comparisons support a scientific claim?}
\sources{Executed classroom experiment: \texttt{Classroom\_Labs.ipynb}; source: \texttt{src/classroom\_labs.py}. Teaching assumptions are explicit.}
\note{174--176 min. Give one minute per pair, then ask for a concise example. Do not accept merely reciting agent names; require MPC, SAC, MILP, WLS, QP or a numerical model where appropriate. The physical simulation stage does not execute field actions. Spoken bridge: Which comparisons support a scientific claim?}
\end{frame}

\begin{frame}[t]{Three checks for scientific claims about agentic control}

\textbf{1. Is the comparison matched?}\par
Same observations, model coverage, limits, tools, candidate budget and test scenarios?\par\medskip
\textbf{2. Which component caused the difference?}\par
Better optimizer, broader model, protected information, language interpretation or coordination?\par\medskip
\textbf{3. What was actually measured?}\par
Constraint compliance, objective, task completion, claim support, latency or human performance?\par\medskip
\uncover<2->{\textbf{Today's result:} the framework makes these questions executable. Extra language models and learned selection did not establish general superiority in the recorded pilot.}

\slidetransition{What single assumption would you change in the next experiment?}
\sources{\href{https://arxiv.org/abs/2606.18789}{PowerAgentBench-SS (2026)}; Actual recorded voltage pilot: \texttt{LLM\_Response\_Selection.ipynb}, E1--E8; raw calls and numerical records in \texttt{results/ensemble\_voltage\_pilot/}.}
\note{176--178 min. This is the scientific takeaway, not a marketing conclusion. Participants should be able to reject an unfair DRL-versus-agentic comparison where the latter receives a better model or more trials. Ask which ablation they would run next. Spoken bridge: What single assumption would you change in the next experiment?}
\end{frame}

\begin{frame}[t]{Your next experiment: change one assumption and keep an audit trail}

\begin{tabular}{@{}>{\raggedright\arraybackslash}p{.20\linewidth}>{\raggedright\arraybackslash}p{.43\linewidth}>{\raggedright\arraybackslash}p{.29\linewidth}@{}}
\toprule Extension & Change one thing & Keep fixed \\
\midrule
Voltage & Forecast uncertainty or a current margin & Controls, model scope and trial budget. \\
Restoration & Priority weights or one tie rating & Fault, graph and baseline constraints. \\
Cyber & Sensor trust or noise assumptions & Physical truth and scoring protocol. \\
Agent design & One role, prompt or selector feature & Tools, test split and numerical gates. \\
\bottomrule
\end{tabular}\par
\par\medskip
\textbf{Deliverable:} one plot, a baseline comparison, a trace explaining the result and a statement of what remains unproven.
\par\medskip
\textbf{After class:} the appendix contains derivations, full code excerpts, recorded responses and training details. It is reference material outside the 180-minute route.

\slidetransition{Choose an extension, then use the appendix when you need its derivation.}
\sources{Executed classroom experiment: \texttt{Classroom\_Labs.ipynb}; source: \texttt{src/classroom\_labs.py}. Teaching assumptions are explicit. Reading guide: \texttt{appendix/Agentic\_AI\_Research\_2026.md}.}
\note{178--180 min. End with an achievable research task, not additional lecture content. Invite participants to choose one extension and write a hypothesis. The teaching session ends here; the appendix is not part of the time budget. Answer final questions and point to the source package and classroom notebook. Spoken bridge: Choose an extension, then use the appendix when you need its derivation.}
\end{frame}

